\chapter{Annexes techniques}

\section{Arborescence simplifiée du backend}
\begin{lstlisting}
backend/app/
  routers/       endpoints HTTP (projects, interview, knowledge-graph,
                 requirements, security, review, documents, ai-settings)
  services/      logique métier (interview, knowledge_graph, requirement,
                 security_engine, review_engine, document_generator,
                 document_rendering, llm_client, ai_settings)
  repositories/  accès aux données (une classe par entité)
  models/        entités SQLAlchemy
  schemas/       schémas Pydantic (validation / sérialisation API)
  domain/        connaissances métier statiques (catalogue de concepts,
                 types d'exigences, référentiels de sécurité, savoir FOCP)
\end{lstlisting}

\section{Exemple d'exigence générée}
Extrait illustratif de la structure d'une exigence générée par SRS (les valeurs sont
indicatives)~:

\begin{lstlisting}
{
  "requirement_key": "SEC-VAL-001",
  "requirement_type": "security",
  "title": "Validation des donnees par un role autorise",
  "priority": "haute",
  "source_concept_key": "workflows.monthly_validation",
  "acceptance_criteria": [
    "Seuls les roles autorises peuvent valider une soumission mensuelle."
  ]
}
\end{lstlisting}

\section{Exemple de diagramme MVP (Mermaid)}
Extrait illustratif du diagramme de conception MVP généré dynamiquement (structure
simplifiée — le contenu réel dépend des exigences confirmées pour chaque projet)~:

\begin{lstlisting}
flowchart TD
    A0["Administrateur"] --> FE["Frontend"]
    A1["Coordinateur regional"] --> FE
    FE --> API["Backend / API"]
    API --> M0["Exigences fonctionnelles"]
    M0 --> DB[(Base de donnees)]
    API --> SEC["Controles de securite (RBAC, audit, chiffrement)"]
\end{lstlisting}

\TODO{Remplacer ces extraits par des exemples réels générés pour le projet Fondation
OCP, une fois un entretien complet mené et un cahier des charges effectivement exporté.}
